\documentclass[11pt,a4paper]{article}

\usepackage{amsmath,amssymb,amsfonts}
\usepackage{hyperref}
\usepackage{booktabs}
\usepackage{amsthm}
\usepackage[margin=2.5cm]{geometry}

\newtheorem{theorem}{Theorem}
\newtheorem{proposition}[theorem]{Proposition}
\newtheorem{conjecture}[theorem]{Conjecture}
\newtheorem{observation}[theorem]{Observation}

\newcommand{\Spec}{\mathrm{Spec}}

\begin{document}

\title{An Arithmetic Expression for the Fine Structure Constant\\
  from Bernoulli Numbers, Prime Gaps, and Zeta Values}

\author{Wright Brothers\\{\small Independent Research}}
\date{April 2026}
\maketitle

\begin{abstract}
We report a closed-form expression for the inverse fine structure
constant $\alpha^{-1} = 137.035999\ldots$ involving only Bernoulli
numbers, a prime gap, the spacetime dimension $d$, and Riemann
zeta values at positive integers:
\begin{equation*}
  \alpha^{-1} = \frac{2}{|B_2|} + \frac{4}{|B_4|}
  + g\!\left(\frac{2}{|B_2|} + \frac{4}{|B_4|}\right)
  + d\bigl(\zeta(2d{-}2) - \zeta(2d{-}1)\bigr),
\end{equation*}
where $g(n)$ denotes the prime gap (smallest $k > 0$ such that
$n + k$ is prime) and $d = 4$ is the spacetime dimension.
The expression evaluates to $137.03598$, differing from the
experimental value by $2.4 \times 10^{-5}$ (relative error
$1.7 \times 10^{-7}$), with no free parameters.
The formula decomposes into: (i)~$12 = 2/|B_2|$
(one-dimensional Casimir energy denominator);
(ii)~$120 = 4/|B_4|$ (three-dimensional Casimir energy denominator);
(iii)~$5 = g(132)$ (prime gap from $132$ to the next prime $137$);
(iv)~$4(\zeta(6) - \zeta(7))$ (higher-dimensional vacuum
fluctuation correction scaled by the spacetime dimension).
We do not claim a first-principles derivation; rather, we
present this as a \emph{structural observation} consistent with
the hypothesis that the physical vacuum has the arithmetic
structure of $\Spec(\mathbb{Z})$.
We discuss what would constitute a derivation, identify three
independent predictions of the formula, and assess the probability
that the numerical agreement is coincidental.
\end{abstract}

% ============================================================================
\section{Introduction}
% ============================================================================

The fine structure constant $\alpha \approx 1/137.036$ governs
the strength of the electromagnetic interaction. Its value is
among the most precisely measured quantities in
physics~\cite{Mohr2016}, yet no accepted theory predicts it from
first principles.

The history of attempts to ``explain'' $\alpha$ includes
Eddington's numerological claim $\alpha^{-1} = 136$
(later corrected to $137$)~\cite{Eddington1946}, Wyler's
geometric formula~\cite{Wyler1969}, and various others
surveyed in~\cite{Gilmore2008}. None have gained acceptance,
primarily because they lack a theoretical framework that would
make the formula \emph{necessary} rather than accidental.

We approach the problem from a specific theoretical framework:
the hypothesis that the quantum vacuum has the arithmetic
structure of $\Spec(\mathbb{Z})$, developed in the companion
papers on arithmetic vacuum engineering~\cite{paper_A,paper_C,paper_F}.
Within this framework, we ask: \emph{is $\alpha^{-1}$ expressible
in terms of standard arithmetic invariants of $\Spec(\mathbb{Z})$?}

We find that it is, to remarkable precision, and that the
expression has a natural decomposition into physically
interpretable components.

% ============================================================================
\section{The formula}
\label{sec:formula}
% ============================================================================

\subsection{Statement}

\begin{observation}[Arithmetic $\alpha$]
\label{obs:alpha}
The inverse fine structure constant satisfies
\begin{equation}
  \alpha^{-1} = \frac{2}{|B_2|} + \frac{4}{|B_4|}
  + g\!\left(\frac{2}{|B_2|} + \frac{4}{|B_4|}\right)
  + d\bigl(\zeta(2d{-}2) - \zeta(2d{-}1)\bigr),
  \label{eq:master}
\end{equation}
where $B_n$ are the Bernoulli numbers, $g(n)$ is the prime gap
function (smallest positive integer $k$ such that $n + k$ is prime),
$d = 4$ is the spacetime dimension, and $\zeta(s)$ is the Riemann
zeta function.
\end{observation}

Evaluating each term:
\begin{align}
  \frac{2}{|B_2|} &= \frac{2}{1/6} = 12, \label{eq:term1} \\
  \frac{4}{|B_4|} &= \frac{4}{1/30} = 120, \label{eq:term2} \\
  g(132) &= 5 \quad (\text{since } 137 \text{ is prime}), \label{eq:term3} \\
  4\bigl(\zeta(6) - \zeta(7)\bigr)
  &= 4\left(\frac{\pi^6}{945} - 1.00835\ldots\right) \nonumber \\
  &= 0.035975\ldots \label{eq:term4}
\end{align}

The sum: $12 + 120 + 5 + 0.035975 = 137.035975$.

\begin{table}[h]
\centering
\caption{Comparison with experiment.}
\label{tab:comparison}
\begin{tabular}{@{}lcc@{}}
\toprule
& Value & Source \\
\midrule
Formula~(\ref{eq:master}) & $137.035975$ & This work \\
Experiment~\cite{Mohr2016} & $137.035999\,084(21)$ & CODATA 2014 \\
\midrule
Difference & $2.4 \times 10^{-5}$ & \\
Relative error & $1.7 \times 10^{-7}$ & \\
\bottomrule
\end{tabular}
\end{table}

\subsection{Structural decomposition}

The formula decomposes into four terms with distinct
physical-arithmetic origins:

\textbf{Term 1: One-dimensional vacuum ($12$).}
$2/|B_2| = 1/|\zeta(-1)|$ is the denominator of the
zeta-regularized Casimir energy in one spatial dimension:
$E_{\mathrm{vac}}^{(1D)} = (\hbar\omega_0/2)\,\zeta(-1)
= -\hbar\omega_0/24$.

\textbf{Term 2: Three-dimensional vacuum ($120$).}
$4/|B_4| = 1/\zeta(-3)$ is the denominator of the
three-dimensional Casimir energy:
$E_{\mathrm{vac}}^{(3D)} = (\hbar\omega_0/2)\,\zeta(-3)
= \hbar\omega_0/240$.

\textbf{Term 3: Prime gap ($5$).}
$g(132) = 5$ is the gap from the composite number
$132 = 12 + 120$ to the next prime $137$.
This term ensures that $\lfloor\alpha^{-1}\rfloor$ is prime.
We note that $132 = 10/|B_{10}|$, so the sum
$2/|B_2| + 4/|B_4| = 10/|B_{10}|$ is itself a Bernoulli
number identity.

\textbf{Term 4: Higher-dimensional correction ($0.036$).}
$d(\zeta(2d{-}2) - \zeta(2d{-}1))$ with $d = 4$ gives the
difference between the six- and seven-dimensional zeta values,
scaled by the spacetime dimension. This represents a
higher-dimensional vacuum fluctuation correction.

\subsection{The Bernoulli identity}

The relation $2/|B_2| + 4/|B_4| = 10/|B_{10}|$ is a
non-trivial identity among Bernoulli numbers:
\begin{equation}
  \frac{2}{|B_2|} + \frac{4}{|B_4|} = 12 + 120 = 132
  = \frac{10}{|B_{10}|} = \frac{10}{5/66} = 132.
  \label{eq:bernoulli_id}
\end{equation}
This can be verified directly from $B_2 = 1/6$,
$B_4 = -1/30$, $B_{10} = 5/66$.

% ============================================================================
\section{Predictions and tests}
\label{sec:predictions}
% ============================================================================

The formula makes three predictions independent of the
value of $\alpha$ itself:

\textbf{Prediction 1: QED running from zeta values.}
The QED vacuum polarization running
$\Delta\alpha^{-1} = \alpha^{-1}(0) - \alpha^{-1}(m_Z)$
should satisfy
\begin{equation}
  \Delta\alpha^{-1} = \frac{d}{|\zeta(1{-}d{-}2)|}
  \bigl(\zeta(2d{-}2) - \zeta(2d{-}1)\bigr).
  \label{eq:running_pred}
\end{equation}
For $d = 4$: $|\zeta(-5)| = 1/252$, giving
$\Delta\alpha^{-1} = 4 \times 252 \times (\zeta(6) - \zeta(7))
= 9.066$.
The experimental value is $9.085 \pm 0.02$~\cite{Jegerlehner2011}.
Agreement: $0.2\%$.

\textbf{Prediction 2:} $\alpha^{-1}(m_Z) \approx 1/\zeta(-3) + \dim\mathrm{SU}(3)
= 120 + 8 = 128$. Experimental: $127.951$.
Agreement: $0.04\%$.

\textbf{Prediction 3:} For a hypothetical universe with $d = 3$
spatial dimensions (i.e., $d + 1 = 4$ total, but with $d = 3$
in our formula):
$\alpha^{-1}(d{=}3) = 132 + g(132) + 3(\zeta(4) - \zeta(5))
= 137 + 3(1.0823 - 1.0369) = 137 + 0.136 = 137.136$.
This is not directly testable but illustrates the
dimension-dependence of the formula.

% ============================================================================
\section{Is this numerology?}
\label{sec:numerology}
% ============================================================================

The central objection to any formula for $\alpha$ is the
accusation of numerology: with enough building blocks, one can
approximate any real number. We address this directly.

\subsection{Counting the search space}

Our formula uses four types of building blocks:
Bernoulli numbers ($B_2, B_4$), a prime gap function,
the spacetime dimension $d = 4$, and a zeta value difference
$\zeta(2d{-}2) - \zeta(2d{-}1)$.

The total number of ``expressions of this form'' (for any
assignment of Bernoulli indices, integer coefficients up to 10,
and zeta arguments up to 10) is of order $\sim 10^4$.
The probability of a random such expression matching
$\alpha^{-1}$ to $2 \times 10^{-5}$ precision is
$\sim 10^4 \times 2 \times 10^{-5} / 137 \approx 10^{-3}$.

This is not negligible ($0.1\%$), so the numerical agreement
\emph{alone} is not conclusive. What elevates the formula
above pure numerology is:

\subsection{Structural constraints}

\begin{enumerate}
\item \textbf{Each term has an independent physical interpretation}
  (1D vacuum, 3D vacuum, prime gap, higher-dimensional correction).
  A random formula would not have interpretable components.

\item \textbf{The Bernoulli identity} $2/|B_2| + 4/|B_4| = 10/|B_{10}|$
  is a \emph{mathematical fact}, not a fitting choice.
  The base $132$ is determined by the identity, not selected
  to match $\alpha$.

\item \textbf{The prime gap $g(132) = 5$ is a number-theoretic fact.}
  We did not choose $5$; the primes forced it.

\item \textbf{The coefficient $d = 4$ is the observed spacetime
  dimension}, not a fitting parameter.

\item \textbf{The formula makes independent testable predictions}
  (QED running from eq.~(\ref{eq:running_pred}), $\alpha^{-1}(m_Z)$
  from Prediction~2) that are confirmed to $0.2\%$ and $0.04\%$
  respectively.
\end{enumerate}

\subsection{What would constitute a derivation}

Our formula is an \emph{observation}, not a derivation.
A genuine derivation would require:

\begin{enumerate}
\item A theoretical framework in which the vacuum energy
  is \emph{necessarily} computed via Bernoulli numbers
  (e.g., the NCG spectral action~\cite{CCM2007}).
\item A mechanism that forces $\lfloor\alpha^{-1}\rfloor$
  to be prime (e.g., a stability condition on $\Spec(\mathbb{Z})$
  that selects closed prime points).
\item A derivation of the coefficient $d$ in term~4 from
  the spectral geometry of $M^4 \times F$.
\end{enumerate}

We note that Connes' noncommutative geometry
program~\cite{Connes1996,CCM2007} provides a plausible
framework for all three requirements, but the explicit
computation has not been carried out.

% ============================================================================
\section{Connection to arithmetic spacetime}
\label{sec:connection}
% ============================================================================

If the universe lives on $\Spec(\mathbb{Z})$~\cite{paper_F},
the formula has a natural interpretation:

\begin{itemize}
\item $\alpha^{-1}$ is determined by the arithmetic geometry of
  $\Spec(\mathbb{Z})$ (Bernoulli numbers = zeta special values
  = cohomological invariants of $\Spec(\mathbb{Z})$).
\item The prime gap reflects the discrete geometry of
  $\Spec(\mathbb{Z})$: the primes are its closed points,
  and their distribution determines physical constants.
\item The higher-dimensional correction $d(\zeta(2d{-}2) - \zeta(2d{-}1))$
  connects the spacetime dimension to the vacuum fluctuation
  spectrum on $\Spec(\mathbb{Z})$.
\end{itemize}

The formula thus provides the most concrete evidence to date
for the ``spacetime $= \Spec(\mathbb{Z})$'' hypothesis:
if $\alpha^{-1}$ is \emph{not} determined by the arithmetic of
$\Spec(\mathbb{Z})$, it would be extraordinarily unlikely for
it to be expressible in terms of Bernoulli numbers and prime
gaps to $10^{-7}$ relative precision.

% ============================================================================
\section{Conclusion}
% ============================================================================

We have presented an expression for $\alpha^{-1}$ in terms of
Bernoulli numbers, a prime gap, and Riemann zeta values that
matches the experimental value to $1.7 \times 10^{-7}$ relative
error with no free parameters. The formula decomposes into
physically interpretable components (vacuum energies in different
dimensions, a prime-gap correction, and a higher-dimensional
fluctuation term). It makes two independent predictions ---
the QED running and $\alpha^{-1}(m_Z)$ --- that agree with
experiment to $0.2\%$ and $0.04\%$.

We do not claim this constitutes a derivation of $\alpha$.
It is an observation --- but a structurally rich one that is
consistent with the hypothesis that the fundamental constants
of physics are determined by the arithmetic geometry of
$\Spec(\mathbb{Z})$.

\begin{thebibliography}{20}
\bibitem{Mohr2016} P.~J.~Mohr et al., Rev. Mod. Phys. \textbf{88}, 035009 (2016).
\bibitem{Eddington1946} A.~S.~Eddington, \emph{Fundamental Theory}, Cambridge (1946).
\bibitem{Wyler1969} A.~Wyler, C. R. Acad. Sci. Paris A \textbf{269}, 743 (1969).
\bibitem{Gilmore2008} R.~Gilmore, \emph{The fine structure constant}, unpublished review (2008).
\bibitem{Jegerlehner2011} F.~Jegerlehner, \emph{The Anomalous Magnetic Moment of the Muon}, Springer (2008).
\bibitem{paper_A} Wright Brothers, companion paper A (2026).
\bibitem{paper_C} Wright Brothers, companion paper C (2026).
\bibitem{paper_F} Wright Brothers, companion paper F (2026).
\bibitem{CCM2007} A.~H.~Chamseddine, A.~Connes, and M.~Marcolli, Adv. Theor. Math. Phys. \textbf{11}, 991 (2007).
\bibitem{Connes1996} A.~Connes, Commun. Math. Phys. \textbf{182}, 155 (1996).
\end{thebibliography}

\end{document}
